\section{\textit{Policy Mining} pada ABAC}

Sebagaimana yang sudah disinggung, spesifikasi kebijakan ABAC secara manual merupakan proses yang mahal dan sulit diskalakan pada sistem berskala besar.
\textit{Policy mining} hadir sebagai solusi otomatis atas persoalan ini.
\textcite{perez-haro_abac_2024} mendefinisikan \textit{policy mining} sebagai prosedur otomatis untuk menghasilkam aturan akses dengan menambang pola atribut-nilai dari \textit{log} akses, yaitu catatan permintaan akses yang telah diberikan izin atau ditolak oleh sistem.
Secara lebih formal, \textcite{diaz-rodriguez_abac_2025} merumuskan masalah \textit{policy mining} ABAC sebagai berikut: diberikan sebuah \textit{log} akses \textit{L}\label{first:symbol-l}, tujuan \textit{policy mining} ABAC adalah menghasilkan kebijakan ABAC yang secara akurat merepresentasikan pola perizinan yang terdapat pada \textit{L}.

Sebagai ilustrasi, berikut merupakan contoh potongan log akses L pada sistem rekam medis rumah sakit.
Setiap baris adalah satu permintaan akses beserta nilai atribut subjek dan objek, aksi yang diminta, dan keputusan yang tercatat.

\begin{longtable}{|>{\centering\arraybackslash}p{0.5cm}|>{\raggedright\arraybackslash}p{1.8cm}|>{\raggedright\arraybackslash}p{1.8cm}|>{\raggedright\arraybackslash}p{2cm}|>{\raggedright\arraybackslash}p{1.8cm}|>{\raggedright\arraybackslash}p{1.4cm}|>{\raggedright\arraybackslash}p{1.8cm}|}
\caption{Contoh potongan log akses ($\mathcal{L}$)}
\label{tab:contoh-log-akses} \\
\hline
\textbf{\#} & \textbf{peran (subjek)} & \textbf{unit (subjek)} & \textbf{jenis objek} & \textbf{unit (objek)} & \textbf{aksi} & \textbf{keputusan} \\
\hline
\endfirsthead

\multicolumn{7}{c}{{\bfseries \tablename\ \thetable{} -- lanjutan dari halaman sebelumnya}} \\
\hline
\textbf{\#} & \textbf{peran (subjek)} & \textbf{unit (subjek)} & \textbf{jenis objek} & \textbf{unit (objek)} & \textbf{aksi} & \textbf{keputusan} \\
\hline
\endhead

\hline
\multicolumn{7}{r}{Lanjut ke halaman berikutnya} \\
\endfoot

\hline
\endlastfoot

% ==================== DATA BARIS ====================
1 & perawat     & ICU       & rekam\_me-dis & ICU & baca  & permit \\
\hline
2 & dokter      & ICU       & rekam\_me-dis & ICU & baca  & permit \\
\hline
3 & perawat     & UGD       & rekam\_me-dis & UGD & baca  & permit \\
\hline
4 & perawat     & ICU       & rekam\_me-dis & UGD & baca  & deny   \\
\hline
5 & staf\_admin & ICU       & rekam\_me-dis & ICU & baca  & deny   \\
\hline
6 & dokter      & ICU       & resep\_obat  & ICU & tulis & permit \\
\hline
7 & perawat     & ICU       & resep\_obat  & ICU & tulis & deny   \\
\hline

\end{longtable}

Dari pola perizinan pada L, \textit{policy mining} berupaya menemukan kebijakan berupa himpunan aturan ringkas yang memetakan setiap permintaan ke keputusan yang sama seperti tercatat.
Untuk potongan di atas, kebijakan berikut sudah memadai:

\begin{enumerate}[label=\textbf{R\arabic*.}, leftmargin=*]
    \item (baca rekam medis oleh klinisi pada unit yang sama).
    JIKA peran $\in \{\text{perawat}, \text{dokter}\}$ DAN jenis objek $=$ \texttt{rekam\_medis} DAN aksi $=$ \texttt{baca} DAN unit(subjek) $=$ unit(objek) $\rightarrow$ izinkan.

    \item (penulisan resep oleh dokter).
    JIKA peran $=$ \texttt{dokter} DAN jenis objek $=$ \texttt{resep\_obat} DAN aksi $=$ \texttt{tulis} DAN unit(subjek) $=$ unit(objek) $\rightarrow$ izinkan.
\end{enumerate}

Permintaan yang tidak dicakup aturan mana pun ditolak secara baku (\textit{default-deny}).

Kebijakan R1, R2 mengizinkan tepat keempat permintaan ber-\textit{permit} dan menolak seluruh permintaan ber-\textit{deny}.
Karena semua \textit{permit} tercakup dan tak ada \textit{deny} yang keliru diizinkan, kebijakan ini merekonstruksi L secara sempurna.
Inilah yang diukur sebagai \textit{F-score} (rata-rata harmonik presisi dan \textit{recall} aturan terhadap L).
Adapun \textit{Weighted Structural Complexity} (WSC) kebijakan ini dihitung dari banyaknya pasangan atribut-nilai yang menyusun kedua aturan.
Kebijakan dengan WSC lebih rendah lebih ringkas dan lebih mudah dikelola administrator.
Kedua metrik inilah yang dipakai untuk menilai kualitas kebijakan pada Hipotesis H1.

Berbagai pendekatan telah dikembangkan untuk menyelesaikan masalah ini.
\textcite{diaz-rodriguez_abac_2025} merangkum bahwa pendekatan-pendekatan tersebut umumnya berbasis salah satu dari beberapa strategi ekstraksi aturan.
Pendekatannya antara lain adalah \textit{association rule mining} yang dipelopori oleh \textcite{xu_mining_2015} yang membentuk aturan secara progresif dari \textit{log permission} namun kinerjanya menurun pada \textit{log} berukuran besar.
Pendekatan lain yang digunakan seperti pembelajaran mesin, seperti yang dilakukan oleh \textcite{iyer_mining_2018} dengan melakukan pendekatan berbasis algoritma PRISM, atau kombinasi \textit{data mining}-statistik yang dilakukan oleh \textcite{chen_polisma_2020}, dan deteksi klaster dengan mengelompokkan permintaan akses sebelum aturan dibentuk.
\textcite{perez-haro_abac_2024} menyoroti bahwa mayoritas pendekatan tersebut mengandalkan strategi berbasis frekuensi, yaitu menetapkan ambang dukungan minimum untuk membentuk aturan.
Namun pendekatan tersebut menghadapi \textit{trade-off} yang nyata.
Ambang tinggi dapat menyisakan banyak sumber daya tanpa aturan, khususnya yang jarang diakses.
Sedangkan ambang rendah dapat memicu ledakan aturan atau \textit{rule explosion} yang tidak andal.
Sebagai alternatif, mereka mengusulkan pemodelan \textit{log} akses sebagai jaringan afiliasi dan menerapkan analisis \textit{biclique} untuk mengekstraksi aturan tanpa memerlukan ambang frekuensi.

Kualitas kebijakan hasil \textit{policy mining} lazim diukur melalui dua dimensi utama, yakni akurasi dan kompleksitas struktural.
Akurasi biasanya diukur menggunakan \textit{F-score}, yaitu rata-rata harmonik antara presisi dan \textit{recall} aturan yang diekstraksi terhadap \textit{log} akses asli \parencite{diaz-rodriguez_abac_2025}.
Kompleksitas struktural diukur melalui \textit{Weighted Structural Complexity} (WSC)\label{first:wsc}, sebuah metrik yang awalnya dirumuskan oleh \textcite{molloy_critical_1995} untuk \textit{role mining} pada RBAC dan kemudian diadaptasi oleh \textcite{xu_mining_2015} untuk konteks ABAC \parencite{perez-haro_abac_2024}.
WSC pada dasarnya menjumlahkan banyaknya pasangan atribut-nilai yang digunakan untuk mendeskripsikan seluruh aturan dalam suatu kebijakan, sehingga kebijakan dengan WSC rendah lebih ringkas dan lebih mudah dikelola oleh administrator.

Kedua metrik ini, akurasi dan WSC, konsisten dengan metrik kualitas kebijakan yang telah ditetapkan pada Hipotesis 1 untuk mengevaluasi CoDA-EDA terhadap \textit{baseline compact-GA/compact-PSO}.
Adapun pendekatan metaheuristik \textit{population-based} seperti ACO \parencite{siyuan_abac_2025} dan UBK \parencite{li_attribute_2026} yang menjadi \textit{baseline} penelitian ini merupakan varian lain dari strategi ekstraksi aturan berbasis pencarian (\textit{search-based}), berbeda dari strategi \textit{association-rule, machine learning}, klaster, maupun \textit{graph-based} yang sudah diuraikan di atas.